\chapter{Infrastructure as system}
\label{vol09:ch15}

\epigraph{The infrastructure that built the country was built on the assumption that someone would maintain it.}{}

\chapmeta{Half-life: durable. Archetypes: scaling, balance, failure. Exercise target: 25. Project track: physical.}

\noindent
Engineering infrastructure is the shared physical foundation of modern civilisation: roads, bridges, water and sewer networks, power grids, telecommunications, ports, airports, dams, levees, public buildings, schools, hospitals. The discipline of treating infrastructure as a system addresses the asset-management questions that emerge when individual structures are aggregated into networks operated and maintained over decades. The engineering of new infrastructure shares the rigor of any design discipline; the engineering of existing infrastructure brings additional concerns: aging, inspection, prioritisation under budget constraints, replacement vs renewal, coupling between systems, and the political economy that determines whether maintenance is funded. This chapter develops infrastructure as engineered systems, asset management at scale, life-cycle costs and total cost of ownership, renewal vs replace decisions, coupled infrastructures and systemic risk, the political economy of infrastructure, and the characteristic failure mode of deferred maintenance.

\section{Infrastructures as engineered systems}\pathtag{standard}

An infrastructure system has thousands or millions of components, operated and maintained for decades or longer, serving large populations with continuous availability. The system is engineered: each component meets a specification, the network is designed to meet load with margin, the operating procedures preserve function under disturbance. The system is also a service: users depend on it without thinking about its components.

Standard infrastructure categories:

\textbf{Transportation:} roads, rail, ports, airports, bridges, tunnels. The American Society of Civil Engineers' (ASCE) Infrastructure Report Card grades US infrastructure C- (in 2025); investments needed are estimated in trillions over decades.

\textbf{Water and wastewater:} treatment plants, distribution and collection networks, dams, levees. Many US systems are over a century old; lead service lines remain in service in many cities.

\textbf{Energy:} the electrical grid, natural-gas networks, oil pipelines, refineries. Aging infrastructure faces both replacement needs and a transition to renewable sources.

\textbf{Telecommunications:} fibre and wireless networks, data centres, undersea cables. Mostly newer than other categories but with rapid obsolescence cycles.

\textbf{Buildings:} schools, hospitals, government buildings, housing. The lifetime is decades but variable maintenance.

\textbf{Coastal and flood-protection:} levees, dykes, breakwaters, beaches. Aging structures meet rising sea levels and intensified storms.

Each category has its own engineering, its own regulatory framework, its own maintenance discipline, and its own political economy. The cross-cutting concerns (asset management, life-cycle costing, system coupling) are the focus of this chapter.

\begin{archetype}[balance: infrastructure as life-cycle balance]
Infrastructure engineering balances three competing objectives: initial cost (the budget to build), operating cost (the budget to run), and renewal cost (the budget to refurbish, replace, expand). Each category trades off against the others: cheaper construction often produces higher maintenance; deferred maintenance becomes catastrophic failure. The engineering discipline is life-cycle balancing: choose the asset, the operating regime, and the maintenance plan to minimise total cost subject to service-level requirements.
\end{archetype}

\section{Asset management at scale}\pathtag{standard}

Asset management is the systematic approach to managing physical assets across their life cycle. The discipline includes: asset inventory (what assets exist, where, in what condition); condition assessment (how the assets are aging); risk-based prioritisation (where investment is most needed); investment planning (how to allocate the budget); performance monitoring (whether the assets are meeting service levels).

ISO 55000 (the international asset-management standard) provides a framework for organisational asset management: policy, strategy, objectives, plans, implementation, performance evaluation. Major infrastructure operators (utilities, transportation agencies, building owners) increasingly conform to ISO 55000 to provide governance and discipline.

Condition assessment methods vary by asset class: visual inspection (bridges, levees), non-destructive testing (ultrasonic, radiographic), instrumented monitoring (strain gauges, accelerometers, corrosion sensors), operational data (failure history, maintenance records), modelling (deterioration curves, finite-element analysis). The cost of assessment is substantial; the trade-off is between assessment cost and the cost of asset failure or premature replacement.

Risk-based prioritisation combines the probability of failure (from condition assessment) with the consequence of failure (operational impact, safety, environment, regulatory). High-probability, high-consequence assets get attention first. The standard frameworks (risk matrices, ALARP, cost-benefit analysis) provide structure; the engineering judgement is in the inputs (probability estimates, consequence assessment, willingness to invest).

Modern asset-management systems combine GIS (geographic information system) data, asset databases, condition data, work-order systems, and analytics. The integration enables data-driven asset management: identify high-risk assets, plan inspections, schedule maintenance, track performance. The systems are increasingly available (Cityworks, Cartegraph, Maximo, SAP) and increasingly mature.

\section{Lifecycle costs and total-cost-of-ownership}\pathtag{standard}

The life-cycle cost of an asset includes acquisition (design, construction), operation (energy, staffing, consumables), maintenance (routine, corrective, preventive), renewal (refurbishment, replacement), and disposal (decommissioning, demolition). The total cost over the asset's life is typically several times the initial cost; for some assets ($100$-year buildings, $50$-year bridges), the multiplier is much higher.

Net present value (NPV) analysis discounts future costs to present value at a chosen discount rate. The standard formula: NPV $= \sum_t C_t / (1 + r)^t$, where $C_t$ is the cost at year $t$ and $r$ is the discount rate. The discount rate is typically the cost of capital ($5\%$ to $10\%$) for private projects, lower ($3\%$ to $5\%$) for public projects emphasising intergenerational equity.

NPV makes long-term decisions consistent: investments with high upfront cost and low operating cost can be compared to investments with low upfront cost and high operating cost. The discount rate strongly influences the comparison: higher rates favour cheaper construction; lower rates favour durable construction.

Equivalent annual cost (EAC) is an alternative formulation: convert the NPV to an equivalent constant annual payment over the asset's life. EAC simplifies comparison of assets with different lifetimes.

Total cost of ownership (TCO) is a closely related concept used in procurement: the buyer compares offers not just on price but on total cost across the asset's life. TCO analysis discourages cheap-but-expensive-to-operate procurement choices.

The discipline of life-cycle costing includes: identifying all relevant costs (some are easy to omit); estimating future costs (subject to uncertainty); choosing the discount rate (consequential and contested); accounting for residual value (some assets have value at end of life); accounting for risk (failure probability and consequence). The engineering work is the input data; the financial analysis is well-established.

\section{Renewal vs replace decisions}\pathtag{standard}

The renewal-vs-replace decision is recurring in infrastructure asset management. The asset has aged; the operating cost is rising; the failure probability is increasing; the original technology may be obsolete. The options:

\textbf{Continue with current maintenance:} accept the rising cost and risk; defer the decision.

\textbf{Refurbish (renewal):} restore the asset to near-new condition through major rebuild. Lower cost than replacement; extends the asset's life. Works when the asset's basic design is still adequate.

\textbf{Replace with similar:} build a new asset of similar specification. Higher cost than renewal; new asset with new design life. Works when the original specification is still appropriate.

\textbf{Replace with different:} build a new asset of different specification (larger capacity, different technology, better resilience). Most expensive option; suited to fundamental changes in service needs.

\textbf{Decommission:} retire the asset without replacement. Suited to assets no longer needed.

The decision combines economic analysis (life-cycle cost of each option), service-level analysis (which options meet the future service requirements), risk analysis (consequences of failure under each option), and policy considerations (regulatory requirements, equity, strategic priorities). The decision is often more political than technical; the engineer's job is to make the technical inputs transparent.

\begin{estimation}
\textbf{Question.} Estimate the optimal replacement age for a class of assets with rising annual maintenance cost and known new-asset cost.

\textbf{Working.} Annual maintenance cost grows from $M_0$ at year $0$ to $M_0(1 + g \cdot t)$ at year $t$ for linear growth; the asset costs $C$ to replace. Total NPV over $T$ years before replacement: $\sum_{t=1}^T M_0(1 + g \cdot t)/(1+r)^t + C/(1+r)^T$. Differentiate with respect to $T$ and set to zero. For $M_0 = 0.01C$, $g = 0.1/\text{year}$, $r = 0.05$, the optimum is around $T \sim 20$ years; the asset is replaced when the annual maintenance reaches about $\sim 3\%$ of new cost. The economic optimum varies with discount rate and growth pattern; the operational optimum may differ (replacement disruption, capital availability, technological turnover).
\end{estimation}

\section{Coupled infrastructures and systemic risk}\pathtag{standard}

Infrastructure systems are increasingly coupled: power feeds water-pumping; water cooling supports power generation; communication networks coordinate power and water; transportation networks support emergency response across all of these. The coupling creates systemic risk: a failure in one system can cascade through coupled systems.

Examples of coupling:

\textbf{Power-water:} water utilities depend on electricity for pumping and treatment. Extended power outages produce water-supply failures, often within hours. The 2003 Northeast blackout caused water-supply outages in many cities.

\textbf{Power-communications:} communication networks (cellular, fibre) depend on electricity for backup generators or battery; long outages exceed backup duration; communication failure then compounds the recovery.

\textbf{Transport-everything:} essentially all infrastructure depends on roads for material transport during construction, maintenance, and emergency response. Transport failures hinder recovery from other infrastructure failures.

\textbf{Information-payment:} modern infrastructure depends on payment systems (for fuel, supplies, services); payment-system failures (cyberattacks, telecom outages) propagate to infrastructure operations.

\textbf{Cooling-data:} data centres require cooling; cooling requires electricity; data centres host applications that coordinate other infrastructure. The chain is fragile under multi-system disturbance.

Systemic risk analysis combines the individual systems' risk with the coupling. Methods: bow-tie analysis, dynamic systems modelling, scenario analysis, stress testing. The engineering challenge is recognising the couplings, characterising the dependencies, and designing resilience that handles multi-system disturbances.

\section{The political economy of infrastructure}\pathtag{mastery}

Infrastructure decisions are made in a political context. The technical analysis is necessary but not sufficient; the engineer who understands the political economy can shape outcomes effectively.

Standard patterns:

\textbf{Construction over maintenance:} political incentives favour visible new construction over invisible maintenance. New bridges get ribbon-cuttings; rebuilt sewers do not. Maintenance backlogs accumulate.

\textbf{Front-loaded benefits, back-loaded costs:} a politician's tenure is shorter than an infrastructure asset's life. Investments with delayed benefits are harder to justify than investments with immediate benefits.

\textbf{Cost-benefit asymmetry:} the cost of building is concentrated (a few visible jurisdictions and timelines); the benefit is diffuse (many users, many decades). The asymmetric political weight favours small visible benefits over large diffuse ones.

\textbf{User finance vs general finance:} user-funded infrastructure (toll roads, water rates) tightens the connection between use and cost but creates equity concerns. Generally-funded infrastructure (general-fund roads, free emergency services) avoids the connection but enables overuse and underfunding.

\textbf{Privatisation:} private operation of public infrastructure (water, transit, energy) trades off private-sector efficiency against public-interest accountability. The standard cases (water privatisation, prison privatisation, healthcare privatisation) have mixed records.

\textbf{Risk shifting:} private financing of public projects (PPPs, design-build-operate, concessions) shifts risk between public and private actors; the resulting allocation often differs from the optimal economic split.

\textbf{Regulatory capture:} the regulator becomes aligned with the regulated industry rather than the public interest. The pattern is common in mature regulated industries.

The engineer who works in the political economy must understand the dynamics; the engineer who ignores the dynamics has technical analyses that go nowhere.

\section{Failure: aging infrastructure and the cost of deferred maintenance}\pathtag{core}

Infrastructure failures have characteristic patterns.

\textbf{Aging without renewal:} the asset continues in service past its design life without major refurbishment or replacement. Failure becomes more likely; consequences become more severe. The Tay Bridge collapse (1879), the I-35W Mississippi River bridge collapse (Minneapolis, 2007), the Genoa Morandi bridge collapse (2018), and many smaller failures show the pattern.

\textbf{Deferred maintenance:} maintenance is postponed for budget or political reasons; the asset deteriorates faster than planned; the rebuild or replacement cost is much higher than the deferred maintenance cost would have been. Reactive maintenance (responding to failures) replaces preventive maintenance (avoiding failures); the asset's life is shortened and the failure rate is higher.

\textbf{Capacity exceeded:} the asset is used beyond its design capacity for too long; the resulting stress causes premature failure. Highway pavements designed for projected traffic ten years out are operated at twice that traffic for thirty years; the consequences appear as roughness, cracking, structural distress.

\textbf{Specification drift:} the asset's specification was correct for original conditions but the conditions have changed. Sea levels rise relative to levees; population grows relative to water supply; vehicle weights increase relative to bridge design. The asset is now under-specified for current conditions.

\textbf{Coupled failure:} as discussed earlier, coupled infrastructures fail together in cascading patterns. The 2003 Northeast blackout, 2012 India blackout, and many smaller multi-system events illustrate the dynamics.

\textbf{Inspection failures:} the inspection program does not detect the deterioration in time; the failure occurs without warning. The 2018 Genoa Morandi bridge failed under a sudden collapse; investigators found significant deterioration that had not been adequately inspected. Inspection itself is engineering work; quality matters.

\textbf{Climate change:} engineered margins designed for historical conditions are eroded by changing climate. Heat tolerances of pavements, drainage capacities, wind loads on structures, snow loads on roofs all change with the climate. Engineered assets need re-evaluation under updated conditions.

\textbf{Knowledge loss:} the engineers who designed and built the original asset retire; the institutional knowledge of design intent, operating envelope, and maintenance needs is lost. New operators may not understand the asset's constraints. Documentation discipline mitigates this; sufficient documentation discipline is rare.

The cost of deferred maintenance is well-documented. ASCE estimates that every dollar of deferred maintenance creates four dollars of future repair cost; the multiplier varies by asset class and deferral duration. The engineering case for sustained maintenance funding is strong; the political case is often weak; the result is the chronic underinvestment that characterises much of US and global infrastructure.

\begin{principle}[The life-cycle-funding principle]
An infrastructure asset is engineered for its life cycle: design, construction, operation, maintenance, renewal, disposal. The discipline is to fund the full cycle, not just the construction. Underfunded maintenance produces premature failure; premature failure produces emergency replacement at higher cost. The life-cycle funding requirement is part of the engineering decision; an asset that cannot be funded over its life cycle should not be built without explicit acknowledgement of the future cost. The engineer's responsibility is making the life-cycle cost transparent to decision-makers; the decision-makers' responsibility is funding accordingly.
\end{principle}

\section{Project}\pathtag{core}

\begin{project}[Study a real infrastructure asset's life cycle; renewal-vs-replace recommendation]
\textbf{Objective.} Study a real infrastructure asset (a bridge, water main, building, road, dam) you can access information about; analyse its life cycle; produce a renewal-vs-replace recommendation with engineering rationale.

\textbf{Method.} Choose an asset for which information is available (your city's water main records, a local bridge's National Bridge Inventory entry, a building's maintenance records). Collect data: original construction, design life, current age, condition history, maintenance expenditures, recent inspections, operating performance. Analyse the life cycle: original cost, accumulated maintenance, projected remaining life, replacement cost. Compute NPV for several alternative paths (continue current maintenance, refurbish, replace, decommission). Identify the political and operational constraints. Recommend a path with engineering rationale.

\textbf{Deliverables.} The asset description with sources; the data collected; the NPV analysis for the alternatives; the recommendation with rationale; a written discussion of the political and operational constraints.

\textbf{Hazard.} The deliverable should not include specific recommendations for actions on public assets without proper authorisation; the analysis is educational.
\end{project}

\section{Exercises}

\subsection*{Calculation}\pathtag{core}

\begin{exercise}
A water main is $80$ years old with a $100$-year design life. Compute the remaining life under linear-deterioration assumptions.
\end{exercise}

\begin{exercise}
An asset costs $\$1\text{M}$ to construct; annual maintenance is $\$20{,}000$; design life is $50$ years; salvage value is $\$50{,}000$. Compute the NPV at $5\%$ discount rate.
\end{exercise}

\begin{exercise}
A bridge has scour-vulnerability rating "1" on the FHWA $0$-$9$ scale. State the engineering implication and the inspection cycle.
\end{exercise}

\begin{exercise}
An asset has rising maintenance cost: $\$10{,}000$ in year $1$, $\$10{,}000 \times 1.1^{t-1}$ in year $t$. Replacement costs $\$200{,}000$. Compute the optimal replacement year at $5\%$ discount rate.
\end{exercise}

\begin{exercise}
A city has $10{,}000\,\text{km}$ of water main; $20\%$ is over $80$ years old; replacement cost is $\$1\text{M/km}$. Compute the cost to replace the aging portion.
\end{exercise}

\subsection*{Derivation}\pathtag{standard}

\begin{exercise}
Derive the equivalent annual cost formula from the NPV expression.
\end{exercise}

\begin{exercise}
Derive the optimal replacement age for an asset with linear deterioration and constant new-asset cost, assuming infinite horizon.
\end{exercise}

\subsection*{Estimation}\pathtag{standard}

\begin{exercise}
Estimate the deferred-maintenance backlog for a US city of $500{,}000$ population.
\end{exercise}

\begin{exercise}
Estimate the cost difference between replacing a $50$-year-old building with a new one of identical specification versus refurbishing the existing.
\end{exercise}

\begin{exercise}
Estimate the impact of a $1\%$ change in discount rate on the optimal replacement age for a typical infrastructure asset.
\end{exercise}

\subsection*{Simulation}\pathtag{standard}

\begin{exercise}
Simulate a Monte Carlo asset-management problem: stochastic deterioration, uncertain failure consequence, capital-budget constraint. Identify the optimal investment strategy.
\end{exercise}

\begin{exercise}
Simulate the cost-of-deferral: an asset whose maintenance is delayed by $5$ years; compute the additional cost.
\end{exercise}

\subsection*{Design}\pathtag{standard}

\begin{exercise}
Design an asset-management program for a small municipal water utility: inventory, condition assessment, prioritisation, planning, monitoring.
\end{exercise}

\begin{exercise}
Design a public-private partnership structure for a new toll road: risk allocation, payment mechanism, performance standards.
\end{exercise}

\begin{exercise}
Design a coupled-infrastructure resilience analysis for a metropolitan area: power, water, communications, transport.
\end{exercise}

\subsection*{Judgement}\pathtag{core}

\begin{exercise}
A team's asset-management analysis recommends increased investment; the political environment prefers tax cuts. State the engineering response.
\end{exercise}

\begin{exercise}
A team's bridge inspection reveals significant deterioration; the public is told the bridge is safe. State the engineering response.
\end{exercise}

\begin{exercise}
A team is choosing between repairing and replacing an aging water-treatment plant. State the engineering criteria.
\end{exercise}

\begin{exercise}
A team is asked to certify that an aging power-grid asset is fit for service. State the engineering considerations.
\end{exercise}

\begin{exercise}
A team's recommended replacement program is approved but underfunded by half. State the engineering response.
\end{exercise}

\begin{exercise}
A team's coupled-infrastructure analysis identifies a single point of failure across power, water, and communications. State the engineering response.
\end{exercise}

\begin{exercise}
A team is asked to evaluate a proposed privatisation of a municipal water utility. State the engineering considerations including risk transfer, service quality, and the cost-of-capital comparison.
\end{exercise}

\begin{exercise}
A team's climate-resilience study suggests that current flood-protection levees are inadequate for projected sea-level rise. State the engineering recommendation.
\end{exercise}

\begin{exercise}
A team's infrastructure includes assets nearing the end of their design life; the asset class has never experienced operation at this age. State the engineering response.
\end{exercise}

\begin{exercise}
A team is asked to support an infrastructure investment recommendation in a public hearing where the engineering analysis is contested by stakeholders with vested interests. State the engineering posture and the role of transparency, alternative scenarios, and explicit acknowledgement of value judgements.
\end{exercise}
